%------------------------------------------------------------------------------%
\documentclass[fleqn,utf8,aspectratio=169,12pt]{beamer}

\usepackage{gaal}

\title{Posições Relativas de Planos}

%------------------------------------------------------------------------------%
\begin{document}

\addtitle

%------------------------------------------------------------------------------%
\section{Posições Relativas entre Planos}

%------------------------------------------------------------------------------%
\begin{frame}{Posições Relativas entre Dois Planos em $\R^3$}
  Dois planos podem ser
  \pause
  \begin{itemize}[<+->]
      \item coincidentes
      \item paralelos distintos
      \item secantes ou concorrentes
  \end{itemize}
\end{frame}

%------------------------------------------------------------------------------%
\begin{frame}{Posições Relativas entre Dois Planos}
  Considere dois planos
  \[
    \pause
    \alpha\colon a_1x + b_1y + c_1z + d_1 = 0
  \]
  \[
    \pause
    \beta\colon  a_2x + b_2y + c_2z + d_2  =0
  \]

  \pause
  Seus vetores normais são
  \[
    \pause
    \vec{n}_1 = \vetortri{a_1}{b_1}{c_1}
    \pause \qquad\qquad
    \vec{n}_2 = \vetortri{a_2}{b_2}{c_2}
  \]
\end{frame}

%------------------------------------------------------------------------------%
\begin{frame}{Planos Paralelos}
  Os planos são paralelos quando seus vetores normais são paralelos
  \[
    \pause
    \vec{n}_2 = \lambda\vec{n}_1
  \]
  \pause
  Eles podem ser paralelos distintos ou coincidentes

  \pause
  \bigskip
  Se
  \(\;
    \vec{n}_2 = \lambda\vec{n}_1
  \;\)
  e
  \(\;
    d_2 = \lambda d_1
  \;\)
  \pause
  os planos são \emph{coincidentes}

  \pause
  \bigskip
  Se os vetores normais são paralelos,
  \pause
  mas as equações não são proporcionais, \\
  \pause
  os planos são \emph{paralelos distintos}
\end{frame}

%------------------------------------------------------------------------------%
\begin{frame}{Planos Secantes}
  Se
  \(
    \vec{n}_1
  \)
  e
  \(
    \vec{n}_2
  \)
  não são paralelos

  \pause
  os planos são \emph{secantes} (ou \emph{concorrentes})

  \pause
  \bigskip
  Nesse caso, eles possuem uma reta em comum
 \end{frame}

 %------------------------------------------------------------------------------%
 \begin{frame}{Interseção entre os Planos Secantes}
  A direção da reta de interseção é perpendicular aos dois vetores normais

  \pause
  \bigskip
  Vetor diretor da reta comum aos planos pode ser
  \[
    \pause
    \vec{v} = \vec{n}_1\times\vec{n}_2
  \]
\end{frame}

%------------------------------------------------------------------------------%
%------------------------------------------------------------------------------%
\begin{question}
  Determine a posição relativa dos planos
  \[
    \pi_{1}\colon 2x - y + 3z - 4 = 0
  \]
  \[
    \pi_{2}\colon 6x - 3y + 9z - 12 = 0
  \]
\end{question}

%------------------------------------------------------------------------------%
\begin{resolution}{Solução}
  Os vetores normais são
  \[
    \vec{n}_{1} = \vetortri{2}{-1}{3}
    \qquad\qquad
    \vec{n}_{2} = \vetortri{6}{-3}{9}
  \]

  \pause
  Como
  \[
    \vec{n}_{2} = 3\vec{n}_{1}
  \]
  \pause
  os planos são paralelos
\end{resolution}

%------------------------------------------------------------------------------%
\begin{resolution}{Solução}
  \onslide<+->{Os coeficientes da segunda equação são três vezes os da primeira}
  \begin{align*}
    \onslide<+->{6   & =3(2)  \\}
    \onslide<+->{-3  & =3(-1) \\}
    \onslide<+->{9   & =3(3)  \\}
    \onslide<+->{-12 & =3(-4)}
  \end{align*}
  \onslide<+->{as duas equações representam o mesmo conjunto de pontos}

  \onslide<+->{Logo, os planos são \emph{coincidentes}}
\end{resolution}

%------------------------------------------------------------------------------%
\endinput

%------------------------------------------------------------------------------%
\begin{question}
  Determine a posição relativa dos planos
  \[
    \pi_{1}\colon 2x - y + 3z - 4 = 0
  \]
  \[
    \pi_{2}\colon 4x - 2y + 6z + 5 = 0
  \]
\end{question}

%------------------------------------------------------------------------------%
\begin{resolution}{Solução}
  Os vetores normais são
  \[
    \vec{n}_{1} = \vetortri{2}{-1}{3}
    \qquad\qquad
    \vec{n}_{2} = \vetortri{4}{-2}{6}
  \]

  \pause
  Como
  \[
    \vec{n}_{2} = 2\vec{n}_{1}
  \]
  \pause
  os planos são paralelos ou coincidentes
\end{resolution}

%------------------------------------------------------------------------------%
\begin{resolution}{Solução}
  Multiplicarmos a primeira equação por 2,
  \pause
  obtemos
  \[
    4x-2y+6z-8=0
  \]
  \pause
  A segunda equação é
  \[
    4x-2y+6z+5=0
  \]
  \pause
  As constantes são diferentes

  \pause
  Portanto, as equações não representam o mesmo plano

  \pause
  Os planos são \emph{paralelos distintos}
\end{resolution}

%------------------------------------------------------------------------------%
\endinput

%------------------------------------------------------------------------------%
\begin{question}
  Mostre que os planos
  \[
    \pi_{1}\colon x + y + z - 3 = 0
  \]
  \[
    \pi_{2}\colon 2x - y + z - 1 = 0
  \]
  são secantes e determine uma equação paramétrica para a reta de interseção
\end{question}

%------------------------------------------------------------------------------%
\begin{resolution}{Solução}
  \onslide<+->{Os vetores normais são
  \[
    \vec{n}_{1} = \vetortri{1}{1}{1}
    \qquad\qquad
    \vec{n}_{2} = \vetortri{2}{-1}{1}
  \]}
  \onslide<+->{Não são paralelos, pois}
  \begin{align*}
    \onslide<+->{1 & = 1} \\
    \onslide<+->{2 & \neq -1}
  \end{align*}
  \onslide<+->{Os planos são secantes}
\end{resolution}

%------------------------------------------------------------------------------%
\begin{resolution}{Solução}
  \onslide<+->{A direção da reta de interseção é}
  \begin{align*}
    \onslide<+->{\vec{v} & = \vec{n}_{1}\times\vec{n}_{2} \\}
    \onslide<+->{& = \begin{vmatrix}
      \vec{i} & \vec{j} & \vec{k} \\
      1       & 1       & 1       \\
      2       & -1      & 1
    \end{vmatrix}}
    \\
    \onslide<+->{& = (2, 1, -3)}
  \end{align*}
\end{resolution}

%------------------------------------------------------------------------------%
\begin{resolution}{Solução}
\begin{columns}[t]
\column{0.4\textwidth}
  \onslide<+->{Ponto comum aos dois planos}
  \onslide<+->{\[
    \begin{cases}
       x + y + z = 3  \\
      2x - y + z = 1
    \end{cases}
  \]}
  \onslide<+->{Escolhemos $z = 0$}
  \onslide<+->{\[
    \begin{cases}
       x + y = 3  \\
      2x - y = 1
    \end{cases}
  \]}
\column{0.3\textwidth}
  \onslide<+->{Somando as equações}
  \begin{align*}
    \onslide<+->{3x & = 4 \\}
    \onslide<+->{x  & = \frac{4}{3} \\}
    \onslide<+->{y  & = 3-x \\}
    \onslide<+->{   & = 3 - \frac{4}{3}\\}
    \onslide<+->{   & = \frac{5}{3}}
  \end{align*}
\column{0.3\textwidth}
  \onslide<+->{Um ponto da reta é
  \[
    A = \left( \frac{4}{3}, \frac{5}{3}, 0 \right)
  \]}
\end{columns}
\end{resolution}

%------------------------------------------------------------------------------%
\begin{resolution}{Solução}
\begin{columns}[t]
\column{0.5\textwidth}
  \onslide<+->{Equação vetorial da reta de interseção}
  \onslide<+->{\[
    \vetortri{x}{y}{z}
    = \Vetortri{\frac{4}{3}}{\frac{5}{3}}{0}
    + t \vetortri{2}{1}{-3}
  \]}
\column{0.5\textwidth}
  \onslide<+->{A equação paramétrica é}
  \onslide<+->{
  \[
    \begin{cases}
      x = \dfrac{4}{3} + 2t \\[5mm]
      y = \dfrac{5}{3} +  t \\[5mm]
      z = -3t
    \end{cases}
  \]}
\end{columns}
\end{resolution}

%------------------------------------------------------------------------------%
\endinput




%%------------------------------------------------------------------------------%
%\section{Lista Mínima}
%
%%------------------------------------------------------------------------------%
%\begin{frame}{Lista Mínima}
%  \begin{enumerate}
%    \item
%  \end{enumerate}
%
%  \vfill
%  Atenção: A prova é baseada no livro, não nas apresentações
%\end{frame}

%------------------------------------------------------------------------------%
\end{document}
%------------------------------------------------------------------------------%
